% ---
% id: statistics/binomial-box/rem:n-nonmono
% kind: rem
% title: "Non-monotonicity of the exact rule in the round count"
% summary: >-
%   The exact rule's endpoints are integer quantiles on a lattice that refines
%   with n, so its interval can widen when n grows by one, as [[thm:monotonicity-limits]] (v)
%   shows, while the n^(−1/2) width of (iv) stands. A consumer may rely on
%   monotonicity in the budget, not in the round count.
% prominence: standard
% status: accepted
% uses: [thm:monotonicity-limits]
% assumes: []
% refs: []
% provenance: newly-derived
% lean:
% ---

\begin{remark}
\label{rem:n-nonmono}
The endpoints $\kappa^{\pm}$ are integer quantiles of a distribution whose
support lattice refines with $n$, so the attainable tail levels move in jumps
and the interval $[k^{-}/n,k^{+}/n]$ can widen when $n$ increases by one, as
Theorem~\ref{thm:monotonicity-limits}\,(v) exhibits. The asymptotic statement
\eqref{eq:width} is unaffected, the width still falling as $n^{-1/2}$, and
no claim of this module needs monotonicity in $n$. A consumer that assumed
``more rounds, tighter set'' would be assuming something false. The
monotonicity that \emph{is} available, and that a consumer may rely on, is
monotonicity in $\ecomp$, Theorem~\ref{thm:monotonicity-limits}\,(i), with the direction
stated there.
\end{remark}
